% Part: computability
% Chapter: computability-theory
% Section: ce-closed-cup-cap

\documentclass[../../../include/open-logic-section]{subfiles}

\begin{document}

\olfileid{cmp}{thy}{clo}

\olsection[اتحاد المجموعات القابلة للتعداد حسابيًا وتقاطعها]{المجموعات
  القابلة للتعداد حسابيًا منغلقة تحت الاتحاد والتقاطع}

للمجموعات القابلة للتعداد حسابيًا خواص انغلاق،
منها ما تقرره المبرهنة الآتية.

\begin{thm}
متى كانت $A$ و$B$ قابلتين للتعداد حسابيًا، كان كل من
$A\cap B$ و$A\cup B$ كذلك.
\end{thm}

\begin{proof}
نستعمل \olref[eqc]{thm:ce-equiv} لننتقل بين توصيفات المجموعات
القابلة للتعداد حسابيًا، ونورد لذلك براهين مختلفة. وإذا خلت إحدى المجموعتين فالنتيجة ظاهرة؛ فلنفرضهما غير خاليتين في البرهانين الأولين.

أما البرهان الأول، فلنأخذ دالة قابلة للحساب~$f$ تعدد $A$، ودالة
قابلة للحساب~$g$ تعدد $B$، ونضع
\begin{align*}
h(x) & = \umin{y}{(f(y) = x \lor g(y) = x)} \text{، و}\\
j(x) & = \umin{y}{(f((y)_0) = x \land g((y)_1) = x)}.
\end{align*}
عندئذ يكون $A\cup B$ مجال $h$، ويكون $A\cap B$ مجال~$j$.

\begin{explain}
وتفسير ذلك بالحساب: إذا كان لنا إجراءان لتعداد $A$ و
$B$، شبه قررنا انتماء $x$ إلى $A\cup B$ بطلب
$x$ في أي من التعدادين؛ وشبه قررنا انتماء $x$ إلى
$A\cap B$ بطلب $x$ في كليهما معًا.
\end{explain}

وأما البرهان الثاني، فنبقي~$f$ معددة لـ$A$ و~$g$ معددة لـ$B$،
ونضع
\[
k(x) = \begin{cases}
f(x/2) & \text{إذا كان $x$ زوجيًا} \\
g((x-1)/2) & \text{إذا كان $x$ فرديًا.}
\end{cases}
\]
فتعدد $k$ المجموعة $A\cup B$، لأن $k$ تتعاقب فيها قيمتا
التعدادين $f$ و~$g$. وأما $A\cap B$ فتحتاج إلى مزيد نظر: إن
كانت $A\cap B$ خالية، كانت قابلة للتعداد حسابيًا على نحو بديهي. وإلا
فليكن $c$ أي عنصر من $A\cap B$، وعرف $l$ بـ
\[
l(x) = \begin{cases}
f((x)_0) & \text{إذا كان $f((x)_0) = g((x)_1)$} \\
c & \text{في غير ذلك.}
\end{cases}
\]
فالدالة $l$ تستعرض أزواج عناصر تعدادي $f$ و$g$، وتخرج
القيمة المشتركة كلما اتفقتا؛ وإن لم تتفقا أخرجت $c$ حتى لا تقف.

وأما البرهان الأخير، فلنأخذ $A$ \emph{مجالًا} للدالة الجزئية القابلة للحساب $m(x)$،
و$B$ مجالًا للدالة الجزئية القابلة للحساب $n(x)$. فتكون $A\cap B$ مجال الدالة
الجزئية $m(x)+n(x)$.

\begin{explain}
ووجهه أن $A$ تجمع القيم التي تتوقف عندها $m$،
و$B$ تجمع القيم التي تتوقف عندها $n$؛ فلا يكون $A\cap B$ إلا مجموعة القيم
التي يتوقف عليها الإجراءان جميعًا.
\end{explain}

وأما رد $A\cup B$ إلى مجموعة قيم توقف فيحتاج إلى
محاكاة $m$ و$n$ على التوازي. فنأخذ $d$ فهرسًا لـ$m$ و$e$ فهرسًا
لـ$n$، فيكون $m=\cfind{d}$ و$n=\cfind{e}$. عندئذ يكون $A\cup B$ مجال
الدالة
\[
p(x) = \umin{y}{(T(d,x,y) \lor T(e,x,y))}.
\]
\begin{explain}
فالدالة $p$ تطلب، عند الدخل $x$، سجل حساب متوقف لـ$m$ أو
لـ$n$؛ ومتى عثرت على أحدهما توقفت.
\end{explain}
\end{proof}

\end{document}
